\begin{frame}{손으로 한 번: Precision · Recall · F1 계산}
  \[
  \text{Precision} = \frac{TP}{TP+FP} = \frac{18}{18+2}
  = \frac{18}{20} = 0.9
  \]
  \[
  \text{Recall} = \frac{TP}{TP+FN} = \frac{18}{18+7}
  = \frac{18}{25} = 0.72
  \]
  \[
  F1 = \frac{2 \times 0.9 \times 0.72}{0.9+0.72}
  = \frac{1.296}{1.62} = 0.8
  \]
  \begin{itemize}
    \item \textbf{Precision 0.9} = ``양성이라고 예측한 20명 중
      18명이 진짜 양성''.
    \item \textbf{Recall 0.72} = ``실제 양성 25명 중 18명만
      잡아냈다'' (7명 놓침).
  \end{itemize}
  \vskip0.4em
  이 예에서 Precision $>$ Recall 이면, 모델이 ``확신이 설 때만
  양성 판정''하는 쪽으로 \textbf{약간 보수적}이라는 신호 ---
  아래에서 다룰 \kb{임계값 조정}이 바로 이 균형을 바꾸는 손잡이.
\end{frame}
